% Glossary terms from ch18: lines of the form \item[Term] Definition.
\item[Kalman filter] The recursive estimator for a linear-Gaussian state-space model: it keeps a Gaussian belief $\Normal(\hat{\state}_k,\mat{P}_k)$ and alternates a predict step through the motion model and an update step with the measurement; it is the exact posterior and the minimum-mean-square-error estimator, and the best linear estimator when the noise is not Gaussian.
\item[Linear-Gaussian state-space model] $\state_k=\mat{F}\state_{k-1}+\mat{B}\ctrl_k+\vect{w}_k$, $\meas_k=\mat{H}\state_k+\vect{v}_k$ with independent Gaussian noises $\vect{w}_k\sim\Normal(\vect{0},\mat{Q})$ and $\vect{v}_k\sim\Normal(\vect{0},\mat{R})$; $\mat{F}$ is the state-transition matrix and $\mat{H}$ the measurement matrix.
\item[Predict step] $\hat{\state}^-_k=\mat{F}\hat{\state}_{k-1}+\mat{B}\ctrl_k$, $\mat{P}^-_k=\mat{F}\mat{P}_{k-1}\mat{F}\T+\mat{Q}$: the belief pushed through the motion model, with the covariance inflated by the process noise.
\item[Innovation] $\vect{y}_k=\meas_k-\mat{H}\hat{\state}^-_k$, the part of the measurement the prediction did not anticipate; its covariance is $\mat{S}_k=\mat{H}\mat{P}^-_k\mat{H}\T+\mat{R}$, and a consistent filter has zero-mean, white innovations with that covariance.
\item[Kalman gain] $\Kgain_k=\mat{P}^-_k\mat{H}\T\mat{S}_k\inv$, the matrix by which the innovation corrects the predicted mean; it weights prediction and measurement by their precisions and corrects unmeasured state components through their covariance with the measured ones.
\item[Joseph form] The covariance update $(\mat{I}-\Kgain\mat{H})\mat{P}^-(\mat{I}-\Kgain\mat{H})\T+\Kgain\mat{R}\Kgain\T$, valid for any gain and always symmetric positive semidefinite; it equals $(\mat{I}-\Kgain\mat{H})\mat{P}^-$ for the Kalman gain.
\item[Constant-velocity model] The kinematic model with state $(\pos,\vel)$, $\pos_k=\pos_{k-1}+\dt\,\vel_{k-1}$, $\vel_k=\vel_{k-1}$, whose unknown accelerations are treated as process noise.
\item[White-noise acceleration] The process-noise model in which the acceleration is continuous-time white noise of intensity $q$ (m$^2$/s$^3$), giving per axis $\mat{Q}=q\begin{psmallmatrix}\dt^3/3&\dt^2/2\\ \dt^2/2&\dt\end{psmallmatrix}$; $q$ is a rate, independent of the sampling interval.
\item[Normalized innovation squared (NIS)] $d^2_k=\vect{y}_k\T\mat{S}_k\inv\vect{y}_k$, the squared Mahalanobis distance of a measurement from its prediction; chi-square with $m$ degrees of freedom for a consistent filter, used to tune $\mat{Q}$ and to gate outliers.
\item[Validation gate] The ellipsoid $d^2\le\gamma$ around the predicted measurement, with $\gamma$ a chi-square quantile; measurements outside it are treated as missing, and with several tracks the gates decide which measurement may be assigned to which track.
\item[Two-point initialization] Starting a track from two measurements one step apart with $\hat{\state}_0=(\meas_b,(\meas_b-\meas_a)/\dt)$ and the matching covariance $\begin{psmallmatrix}\mat{R}&\mat{R}/\dt\\ \mat{R}/\dt&2\mat{R}/\dt^2\end{psmallmatrix}$.
\item[Observability] The property that the observability matrix $(\mat{H};\mat{H}\mat{F};\dots;\mat{H}\mat{F}^{n-1})$ has full rank, so that the state can be recovered from a finite number of noise-free measurements; without it the covariance of the unobserved components grows without bound.
\item[Prediction covariance] $\mat{P}_{k+j\mid k}=\mat{F}^j\mat{P}_k(\mat{F}^j)\T+\sum_{i<j}\mat{F}^i\mat{Q}(\mat{F}^i)\T$, the uncertainty of the state $j$ steps ahead given the measurements up to $k$; it grows without bound in $j$ and is what the local safety layer inflates an intruder's radius with.
